\documentclass[11pt,a4paper]{article}
\usepackage[margin=23mm]{geometry}
\usepackage{amsmath,amssymb,amsthm,hyperref}
\newtheorem{theorem}{Theorem}
\newtheorem{lemma}[theorem]{Lemma}
\newtheorem{corollary}[theorem]{Corollary}
\newcommand{\Del}{\Delta}
\newcommand{\om}{\omega}
\newcommand{\Om}{\Omega}
\title{Primitive rectangular partitions and planar density}
\author{Human Compression project: working application for F008}
\date{}
\begin{document}\maketitle
This working application uses the complete companion
\emph{Divisor concentration and primitive bases}. Its weighted moment,
exceptional-divisor and multiples lemmas are acquired in full. It recovers
the planar density argument of Lebowitz--Lockard and Tenenbaum, retained
preprint dated 28 August 2026, without asserting a three-dimensional
density theorem. We explicitly retain the Hardy--Ramanujan normal-order
theorem $\om(n),\Om(n)=(1+o(1))\log_2 n$ on a set of density one.
No novelty or compression factor is asserted before distinct review.

\section{Geometry supplies a bounded multiset}
Consider a partition of an integral $k$-cube of side $n$ into $R\ge2$
axis-parallel boxes of equal volume and integral side lengths. Suppose
the box lengths along one specified axis are pairwise distinct. Such a
partition will be called strong. For $k=2$ this is the convention of
the retained source. Scaling preserves the property. A primitive side
length is one divisible by no smaller side length admitting such a
partition.

\begin{lemma}\label{trace}
For a primitive strong partition and a proper divisor $h>1$ of $n$,
there is a coordinate line whose successive box lengths give
\[
 n=\sum_{d\mid n^k}\epsilon_d d,\qquad
 0\le\epsilon_d\le R\le\tau(n^k),\qquad
 \sum_d\epsilon_d\ge2,
\]
with every selected $d$ in $[n/\tau(n^k),n)$ and at least one selected
$d$ not divisible by $h$. When $k=2$, every coefficient is zero or one.
No prime side length admits a strong partition in any dimension.
\end{lemma}
\begin{proof}
Let $V=n^k/R$ be the common volume. Each side length divides $V$ and
therefore $n^k$. The distinct lengths on the specified axis show
$R\le\tau(n^k)$. Since all other sides are at most $n$, every side is
at least $V/n^{k-1}=n/R$.

If all side lengths were divisible by $h$, scaling the entire partition
by $1/h$ would contradict primitivity. Choose a box and a side length
not divisible by $h$. A line parallel to that side through the box's
interior, chosen to avoid all other box boundaries, cuts a partition of
an interval of length $n$. The selected segment is shorter than $n$,
because $h\mid n$, so there are at least two segments. Every segment
on this line is consequently shorter than $n$. Group equal lengths to
obtain the displayed coefficients. In dimension two, distinct heights
and equal areas imply distinct widths, so either coordinate direction
has no repeated segment length. In higher dimensions we retain the
bounded coefficients; distinct heights alone do not imply distinct
widths or depths.

If $n=p$ is prime, all side lengths are in $\{1,p\}$. Some box has
length one on the distinguished axis: otherwise all its lengths would
be $p$, contradicting pairwise distinctness and $R\ge2$. A line in
that direction through this box must consist entirely of unit segments.
It meets at least two boxes with equal distinguished length, a
contradiction.
\end{proof}

\begin{lemma}[Few sums, few prime factors]\label{count}
Fix $k,C>0$ and $0<c<1$. Let $x$ be large and
$z=x^{1/(4\log_2x)}$. Suppose a collection of integers $b\le x$ has
representations $b=mp$, with $p>z$ prime and $p\nmid m$, such that
\begin{enumerate}
\item $\tau(m^k)\le(\log x)^C$;
\item there is a nonzero integer $D\le x$, divisible by $p$, that is
      a sum of at most $(\log x)^c$ distinct divisors of $m^k$, each
      with coefficient between one and $(\log x)^C$.
\end{enumerate}
Then the collection has size $O_A(x/(\log x)^A)$ for every fixed $A$.
\end{lemma}
\begin{proof}
For a fixed $m\le x/z$, put $J=\tau(m^k)$ and
$K=\lfloor(\log x)^c\rfloor$. The possible support and coefficient
lists, including shorter lists, number at most
\[
 (K+1)(J+1)^K(1+(\log x)^C)^K
 =\exp\{O((\log x)^c\log_2x)\}.
\]
We may overcount orderings and lists giving the same $D$. Each positive
$D\le x$ has at most $\log x/\log2$ distinct prime divisors. Thus the
number of pairs $(m,p)$ is at most
\[
 \frac{x}{z}\exp\{O((\log x)^c\log_2x)\}\log x.
\]
The negative exponent $-\log x/(4\log_2x)$ dominates both the positive
term and $A\log_2x$, since $c<1$.
\end{proof}

\section{The planar primitive basis}
Let $B$ be the primitive basis of the set of planar Mondrian side
lengths, under the strong convention just stated. Fix
$\varepsilon=1/100$ and split $B=B_1\cup B_2$ according as
$\Om(b)>(1+\varepsilon)\log_2b$ or not.

\begin{theorem}\label{planar}
The set of planar Mondrian side lengths has an asymptotic density
strictly less than one.
\end{theorem}
\begin{proof}
The companion's exceptional-divisor lemma makes the upper densities of
the multiples of $B_1\cap[N,\infty)$ tend to zero. We prove
\[
 \#(B_2\cap[1,x])\ll x/(\log x)^{11/10}.
\]
The part $b\le\sqrt x$ is harmless. Put $L=\log x$, $\xi=\log_2x$,
$z=x^{1/(4\xi)}$. For $b>\sqrt x$ in $B_2$, its largest prime factor
$p$ exceeds $z$: otherwise
$b\le z^{\Om(b)}\le x^{(1+\varepsilon)/4}<\sqrt x$.
The integers divisible by $p^2$ for some prime $p>z$ number at most
$x\sum_{j>z}j^{-2}=O(x/z)$. Discard them and write $b=mp$, with
$p\nmid m$ and $m\le x/z$.

We have
\[
 R\le\tau(b^2)\le3^{\Om(b)}
 \le(\log b)^{(1+\varepsilon)\log3}< (\log b)^2
\]
for sufficiently large $b$. All segment lengths in Lemma~\ref{trace}
therefore belong to
$I_b=[b/(\log b)^2,b)$.
The number $T(b)$ of divisors of $b^2$ in $I_b$ satisfies
$T(b)\ll\xi\Del(b^2)$, by covering that interval with unit
logarithmic intervals.

The companion moment theorem with $r=2,t=2,y=1/4$ gives
\[
 \sum_{b\le x}\Del(b^2)^2(1/4)^{\om(b)}
 \ll xL^{-3/4+o(1)}.
\]
Since $\om(b)\le\Om(b)\le(1+\varepsilon)\log_2b$ on $B_2$,
\[
 \sum_{\substack{b\le x\\b\in B_2}}\Del(b^2)^2
 \ll xL^{(1+\varepsilon)\log4-3/4+o(1)}.
\]
Markov's inequality shows that $T(b)>L^{9/10}$ holds for at most
\[
 xL^{(101/100)\log4-3/4-9/5+o(1)}
 =O(x/L^{11/10})
\]
integers in $B_2$, because
$9/5+3/4-(101/100)\log4>11/10$.

For each remaining $b$, the prime $p$ is a proper divisor by
Lemma~\ref{trace}. Apply that lemma with $h=p$ and let $D$ be the sum
of its terms not divisible by $p$. Such a term exists; all terms are
positive. Reduction modulo $p$ gives $p\mid D$, and $1\le D\le b$.
Each selected divisor not divisible by $p$ divides $m^2$. There are at
most $T(b)\le L^{9/10}$ different such divisors, and
$\tau(m^2)\le\tau(b^2)\le L^2$. Lemma~\ref{count}, even allowing
coefficients up to $L^2$, bounds these remaining $b$ by $O(x/L^2)$.
This proves the displayed count.

Partial summation gives $\sum_{b\in B_2}1/b<\infty$. Its tails thus
have multiples of upper density tending to zero. The same was proved
for $B_1$, hence for $B$. The companion's density lemma applies, since
$1\notin B$. Scaling and the existence of a least positive divisor
with a partition show that the full set of Mondrian side lengths is
exactly the set of multiples of $B$, proving the theorem.
\end{proof}

\section{Normal orders and the number of rectangles}
\begin{lemma}
For every fixed integer $k\ge1$,
$\tau(n^k)=(\log n)^{\log(k+1)+o(1)}$ on a set of density one.
\end{lemma}
\begin{proof}
At a prime exponent $\nu\ge1$,
$k+1\le k\nu+1\le(k+1)^\nu$. Multiplication gives
$(k+1)^{\om(n)}\le\tau(n^k)\le(k+1)^{\Om(n)}$.
The retained Hardy--Ramanujan theorem gives both bounds with the
claimed exponent.
\end{proof}

\begin{lemma}\label{normal}
For fixed $k\ge2$,
$\Del(n^k)=(\log n)^{\log(k+1)-1+o(1)}$ on a set of density one.
\end{lemma}
\begin{proof}
The logarithms of the divisors lie in $[0,k\log n]$. At most
$k\log n+2$ unit intervals cover them, so
$\Del(n^k)\ge\tau(n^k)/(k\log n+2)$, proving the lower bound.

For the upper bound fix $\eta>0$ and choose a sufficiently small
fixed $t>0$ that
\[
 \frac{(k+1)^t-1-t}{t}<\log(k+1)-1+\eta/2.
\]
The companion moment theorem with $y=1$ applies in its explicitly
proved fractional range. Markov's inequality bounds the proportion of
$n\le x$ with
$\Del(n^k)>(\log x)^{\log(k+1)-1+\eta}$ by a quantity tending to zero.
Replacing $x$ by $n$ for $n>\sqrt x$ changes only fixed constants.
The part below $\sqrt x$ has density zero. Applying the result for
$\eta=1,1/2,1/3,\ldots$ and choosing increasing thresholds gives
the asserted $o(1)$ statement on a density-one set.
\end{proof}

\begin{corollary}
Every planar Mondrian decomposition has
$R(n)\le(\log n)^{\log3-1+o(1)}$ on a set of density one.
\end{corollary}
\begin{proof}
First $R\le\tau(n^2)=(\log n)^{\log3+o(1)}$. Every height lies
between $n/R$ and $n$ and divides $n^2$. Cover that interval by
$\lceil\log R\rceil+1=O(\log_2n)$ unit logarithmic intervals.
Distinctness of heights gives
$R\ll\Del(n^2)\log_2n$; apply Lemma~\ref{normal}.
\end{proof}

\section{The unresolved three-dimensional argument}
Lemma~\ref{trace} shows what survives in higher dimension. Choosing a
side not divisible by a proper divisor may require a different axis
from the one with distinct heights. Its lengths must be treated as a
multiset. Lemma~\ref{count} accommodates those bounded multiplicities.
The missing ingredient is a suitable exceptional-set bound for the
number of divisors of $m^3$ near $mp$. No such bound is established here.

There is also a literal Fourier error in the retained source, physical
page 11. With $\tau(v,\theta)=\sum_{d\mid v}d^{i\theta}$ and
$u=\theta\log p$, the correct identity is
\[
 |\tau(p^3,\theta)|^2
 =|1+e^{iu}+e^{2iu}+e^{3iu}|^2
 =4+6\cos u+4\cos(2u)+2\cos(3u).
\]
It follows by counting ordered pairs of exponents with each difference.
At $u=\pi/2$ the value is zero, whereas the source's
$16\cos^6(u/2)$ is two. Correcting this identity alone does not
justify its later prime-pair estimates, so this note does not certify
Theorem 1.3. It neither disproves that theorem nor answers the Mondrian
nonexistence conjecture. The multiplicative-weight and improved-error
remarks of Section 4 also remain outside the present proof.
\end{document}
